% Nouveau EGA V (cote 25): archive cote 25, site page 104 (PDF page 105).
% typescript. Draft compared: none.
% AI-transcribed by Claude Opus 5.5 (claude-opus-5-5, adaptive thinking, effort high), Message Batches API, batch ega25b (msgbatch_01A7n5YuXubUPNK1zsepi8RG), 2026-11-13; not yet checked by eye.
\scanpage{25/104}
\mspage{7}

\uline{mes} \uline{constants} rel. à S. 2) n=1, alors il faut prendre les homomorphis-%
mes \xout{\unsure{partout surjec}} $\underline{E}'\to \underline{E}\otimes\underline{M}$ qui sont \uline{surjectifs}, comme on vérifie
aussitôt, et l'homomorphisme correspondant \unsure{$g$}: $P\to P'$ n'est autre que le
composé $P(\underline{E})\xrightarrow{\sim}P(\underline{E}\otimes\underline{M})\to P(\underline{E}')$ , où le premier hom est l'isomorphisme
canonique envisagé dans le Chap II, et le d euxième l'immersion fermée
canonique déduite de l'épimorphisme $\underline{E}'\to\underline{E}\otimes\underline{M}$ . Si on appelle \uline{linéaires}
les hom de P dans P' qui se décrivent comme un tel composé, on voit donc
que les morphismes $g:P\to P'$ qui sont de degré 1 sont exactement ceux qui
sont \uline{linéaires}.

Pour terminer, déterminons les isomorphismes de P avec P':
\uline{Théorème} 1.6. Soient S un préschéma, $P=P(\underline{E})$ et $P'=P(\underline{E}')$ deux fibrés
projectifs sur S\inserted{\unsure{définis par} $\underline{E}$, $\underline{E}'$ \unsure{loc. libres de type fini},}. Alors tout S-isomorphisme $g$: $P\xrightarrow{\sim}P'$ est définissable
comme le composé $P(\underline{E})\xrightarrow{\sim}P(\underline{E}\otimes\underline{M})\to \underline{P}(\underline{E}')$, où\inserted{$\underline{M}$ est un Module inversible sur S,} la premier hom est l'isom
can de Chap II , et le \unsure{deuxième|dauxième} est l'isomorphisme \xout{canonique} déduit
d'un isomorphisme u: $E'\xrightarrow{\sim}\underline{E}\otimes\underline{M}$ . Lorsque les fibres de \unsure{P} sont non
vides (resp. de dim $\geq 1$) alors $\underline{M}$ (resp. le couple $(\underline{M},u)$) est déterminé
à isomorphisme unique près en termes de g.

D'après les considérations précédentes, on est ramené à prouver
que g est de degré 1, ce qui nous ramène au cas où S est le spectre d'un
corps,\inserted[typed]{et bien sur on peut supposer dim $P\geq 1$.} Mais notons que $\underline{O}_P(1)$ est alors caractérisé intrinséquement (i.e.
indépendamment de la façon dont P a été réalisé comme fibré projectif)
comm\xout{ex}e le générateur de Pic(P) (parmi les deux générateurs $\underline{O}_P(1)$ et
$\underline{O}_{\unsure{P}}(-1)$ ) qui est \uline{ample} ; par suite si g : $P\to P'$ est un isomorphisme,
alors $g^{*}(\underline{O}_{P'}(1))$ est isomorphe à $\underline{O}_P(1)$ , et on gagne. - Sous forme local
% NOTE: the right edge of the scan cuts off the end of "local[e]"
et moins savante, on peut énoncer:
\uline{Corollaire} 1.7. Sous les conditions de 1.6. , tout\inserted{S-}isomorphisme $g:P\to P'$
peut se décrire, au voisinage de tout $s\in S$, à l'aide d'un isomorphisme
u: $\underline{E}|U \simeq E'|U$ , ce dernier étant bien déterminé mod mult par un élément
\margin{104}
